% !PRINTED 547
\weberpage{547}
%<<CONT
Gleichung, deren Wurzeln die zum zweiten Geschlecht gehörigen Klasseninvarianten sind, und die daher konjugiert imaginär sind, während die Wurzeln von (5) reell sind. Daraus ergibt sich die Größenbestimmung:
\[
(a+b\sqrt p)^2 > 4 > (a-b\sqrt p)^2;
\]
also müssen $a$ und $b$ gleiches Zeichen haben. Da aber $x$ nach (4) positiv ist, so müssen $a$ und $b$ beide positiv sein.

Um $a$ und $b$ wirklich zu finden, braucht man nur den Ausdruck auf der linken Seite von (5) nach (4) auf wenige Dezimalen zu berechnen, wobei es weitaus hinreichend ist,
\[
f(\sqrt{-m}) = f_1(\sqrt{-m}) = e^{\frac{\pi\sqrt m}{24}}
\]
zu setzen, und die so gefundenen Dezimalen mit den Dezimalen von $\sqrt p$ zu vergleichen, um alsbald $b$ und sodann $a$ zu erhalten. Die Rechnung ist überaus einfach und gibt folgende Resultate:
\[
\begin{aligned}
m &= 14, & x + \frac1x &= 1 + \sqrt2,\\
m &= 17, & x + \frac1x &= \frac{1+\sqrt{17}}{2},\\
m &= 34, & x + \frac1x &= \frac{3+\sqrt{17}}{2},\\
m &= 46, & x + \frac1x &= 3 + \sqrt2,\\
m &= 49, & x + \frac1x &= 2 + \sqrt7,\\
m &= 73, & x + \frac1x &= \frac{5+\sqrt{73}}{2},\\
m &= 82, & x + \frac1x &= \frac{15+\sqrt{41}}{2},\\
m &= 97, & x + \frac1x &= \frac{9+\sqrt{97}}{2},\\
m &= 142, & x + \frac1x &= 9 + 5\sqrt2,\\
m &= 193, & x + \frac1x &= 13 + \sqrt{193}.
\end{aligned}
\]

\begin{center}\rule{3cm}{0.4pt}\end{center}
